\emph{完备性定理}是\emph{可靠性定理}的逆命题。
它的一种形式是：若 $\Gamma \Entails !A$ 则 $\Gamma \Proves !A$；另一种形式是：若 $\Gamma$ 是一致的，则它是可满足的。
我们证明了第二种形式（并由后者推出前者）。
该证明较为复杂，分为若干步骤。
我们从一个一致的集合~$\Gamma$ 出发。
首先，我们添加无穷多个新的常项符号~$c_i$，以及形如 $\lexists[x][!A(x)] \lif !A(c)$ 的公式，其中扩充语言中每个含自由变元的公式~$!A(x)$ 都与某个新常项配对。
由此得到一个包含~$\Gamma$ 的\emph{饱和}一致语句集，且该集合仍然是一致的。
接着，我们将该集合扩充为一个\emph{完全一致集}。
完全一致集具有如下良好性质：对任意语句~$!A$，要么 $!A$ 要么 $\lnot !A$ 属于该集合（但不会两者皆属于该集合）。
由于我们从饱和集出发，现在便得到了一个包含~$\Gamma$ 的饱和、完全、一致的语句集~$\Gamma^*$。
由此集合，现在可以定义一个结构~$\Struct{M}$，使得 $\Sat{M(\Gamma^*)}{!A}$ 当且仅当 $!A \in \Gamma^*$。
特别地，$\Sat{M(\Gamma^*)}{\Gamma}$，即 $\Gamma$ 是可满足的。
若语言中包含等词 $=$，则构造略为复杂。

完备性定理有两个重要的推论。
\emph{紧致性定理}指出：$\Gamma \Entails !A$ 当且仅当存在某个有限子集 $\Gamma_0 \subseteq \Gamma$ 使得 $\Gamma_0 \Entails !A$。
一个等价的表述是：$\Gamma$ 是可满足的，当且仅当每个有限子集 $\Gamma_0 \subseteq \Gamma$ 都是可满足的。
紧致性定理可用于证明具有特定性质的模型的存在性。
例如，我们可以用它来证明，任何具有任意大有穷模型的理论，都存在无穷模型。
这尤其意味着，\emph{有穷性}在一阶逻辑中是不可表达的。
第二个推论是 \emph{Löwenheim-Skolem}定理，它指出每个可满足的~$\Gamma$ 都有一个可数模型。
这进而表明，\emph{不可数性}在一阶逻辑中也是不可表达的。